\documentclass[conference]{IEEEtran}
\IEEEoverridecommandlockouts

\usepackage{cite}
\usepackage{amsmath,amssymb,amsfonts}
\usepackage{algorithmic}
\usepackage{graphicx}
\usepackage{textcomp}
\usepackage{xcolor}
\usepackage{booktabs}
\usepackage{hyperref}

\def\BibTeX{{\rm B\kern-.05em{\sc i\kern-.025em b}\kern-.08em
    T\kern-.1667em\lower.7ex\hbox{E}\kern-.125emX}}

\begin{document}

\title{Analyzing India's Rising Currency in Circulation: An Explainable Machine Learning Approach to the ``Capable-but-Cash'' Paradox}

\author{\IEEEauthorblockN{Aarush Garg}
\IEEEauthorblockA{\textit{Dept. of Information Science \& Engineering} \\
\textit{RV College of Engineering}\\
Bengaluru, India \\
aarushgarg@rvce.edu.in}
\and
\IEEEauthorblockN{Arya Wadhwa}
\IEEEauthorblockA{\textit{Dept. of Information Science \& Engineering} \\
\textit{RV College of Engineering}\\
Bengaluru, India \\
aryawadhwa@rvce.edu.in}
\and
\IEEEauthorblockN{Dilraj Singh}
\IEEEauthorblockA{\textit{Dept. of Information Science \& Engineering} \\
\textit{RV College of Engineering}\\
Bengaluru, India \\
dilrajsingh@rvce.edu.in}
}

\maketitle

\begin{abstract}
Despite the exponential proliferation of India's Unified Payments Interface (UPI)---which scaled from 16.58 billion transactions worth INR~23.49 lakh crore in October 2024 to 25.47 billion monthly transactions worth INR~31.53 lakh crore in October 2026 \cite{npci2026}---India's physical banknotes in circulation expanded by 11.9\% in FY 2025--26 to INR~41.24 lakh crore as of March 2026, with total Currency in Circulation (CiC) exceeding INR~43 lakh crore later in 2026 \cite{rbi2026}. Pooling official World Bank Global Findex India microdata across the 2021 and 2025 (2024 wave) editions ($N = 6,000$), we isolate the ``Capable-but-Cash'' cohort ($N_{\text{cap}} = 2,172$) and find that 55.7\% (58.4\% in 2021; 54.0\% in 2024) remain cash-only at point of sale. In the 2024 Findex wave (\texttt{fin25e4}), 77.4\% of cash-reliant capable adults cite habit persistence, 9.8\% cite merchant cash-only acceptance, 6.5\% cite distrust, and only 2.3\% cite higher digital costs. Crucially, merchant-side resistance interacts with tax compliance behavior: automated SFT reporting under Section 285BA into the Annual Information Statement (AIS) and GST registration thresholds (INR~20L/INR~40L under Sections 22/63) incentivize informal micro-merchants to suppress digital turnover. We formulate a strictly leak-free empirical architecture combining survey-weighted Spline-Logit, Regularized XGBoost/LightGBM, a Platt-Calibrated 4-Family Stacking Ensemble (5$\times$3 CV ROC-AUC $= 0.8470 \pm 0.022$), Fairlie non-linear decompositions, and TreeSHAP interpretability to guide interventions intercepting the ``Cash-Out Loop'' and harmonizing micro-merchant tax safe-harbors.
\end{abstract}

\begin{IEEEkeywords}
Financial Inclusion, Digital Payments, UPI, Currency Demand Paradox, Tax Compliance Friction, AIS/TIS Tracking, Explainable AI, TreeSHAP, XGBoost, Behavioral Economics, Global Findex.
\end{IEEEkeywords}

\section{Introduction}
The modern Indian retail payment ecosystem is recognized globally as a pioneering benchmark in Digital Public Infrastructure (DPI) \cite{worldbank2023}. Through the coordinated deployment of the ``India Stack''---anchored by Aadhaar biometric identification, universal account access under the Pradhan Mantri Jan Dhan Yojana (PMJDY), and the open, interoperable Unified Payments Interface (UPI) developed by the National Payments Corporation of India (NPCI)---the marginal technical and economic friction of digital retail payments has been drastically reduced \cite{dasilva2019, npci2026}.

Under orthodox monetary theory and classical payment adoption models (e.g., the Baumol-Tobin cash inventory framework), the near-zero marginal cost of electronic transactions should substitute physical currency as a medium of exchange, triggering a structural contraction in banknote demand \cite{baumol1952, tobin1956}. 

However, empirical data from the Reserve Bank of India (RBI) reveals an enduring macroeconomic divergence:
\begin{itemize}
    \item \textbf{The Digital Surge:} Monthly UPI retail transaction volumes scaled from 14.96 billion in August 2024 (16.58 billion worth INR~23.49 lakh crore in October 2024) to 25.47 billion transactions worth INR~31.53 lakh crore in October 2026, totaling INR~314.23 lakh crore in FY 2025--26 \cite{npci2026}.
    \item \textbf{The Currency Demand Paradox:} Concurrently, India's banknotes in circulation grew by 11.9\% YoY from INR~36.87 lakh crore (March 2025) to INR~41.24 lakh crore (March 2026), with total CiC exceeding INR~43 lakh crore later in 2026---concentrated in INR~500 notes (~85\%--86\% of total note value) \cite{rbi2026, awasthy2023}.
\end{itemize}

While central bank economists decouple this aggregate trend into transactional velocity versus precautionary/store-of-value holdings in the informal economy \cite{goodhart2020, awasthy2023}, macroeconomic indicators obscure granular individual-level behavioral decision-making.

Most existing financial inclusion literature assesses inclusion through an \textit{access-centric paradigm}---measuring bank account ownership rates, smartphone penetration, and telecom connectivity \cite{demirguc2022}. While India successfully expanded formal account ownership from 35\% in 2011 to 77.5\% (approx. 78\%) in 2021 \cite{demirguc2022}, microdata demonstrates that universal access does not automatically dictate active digital usage:
\begin{itemize}
    \item \textbf{Account Inactivity Rate:} India exhibits an account inactivity rate of 35.0\% (no deposits or withdrawals in the preceding 12 months), significantly higher than the developing economy average of 13.0\% \cite{demirguc2022}.
    \item \textbf{Cash-Only Merchant Scale:} Over 670 million banked adults in India pay merchants exclusively in cash \cite{demirguc2022}. India alone accounts for more than 40\% of the 1.6 billion banked cash-only consumers globally.
\end{itemize}

We designate this phenomenon the \textbf{``Capable-but-Cash'' Paradox}. This study addresses the following central research question:
\textit{Among Indian adults who already possess an active formal bank account, a smartphone, and internet access, which behavioral, demographic, and psychological factors are most strongly associated with continued cash reliance for in-store transactions, and how do merchant-side tax frictions interact with consumer choice?}

Rather than relying strictly on classical parametric regressions---which assume linear log-odds and struggle to identify complex feature interactions---we develop an \textbf{Explainable Machine Learning} architecture combining regularized XGBoost \cite{chen2016} with \textbf{TreeSHAP} (Tree-based SHapley Additive exPlanations) \cite{lundberg2020}. This framework allows us to extract an evidence-led, game-theoretically consistent ranking of behavioral barriers, moving beyond binary prediction toward actionable policy recommendations.

\section{Theoretical Framework \& Related Work}

\subsection{The Currency Demand Paradox}
Central banking literature across both advanced and developing economies has documented the ``paradox of banknotes''---wherein aggregate currency demand expands even as electronic alternatives flourish \cite{goodhart2020, arango2015, agarwal2023}. In India, Reserve Bank of India empirical working papers \cite{awasthy2023} and macroeconomic evaluations \cite{kumar2025, minj2025} demonstrate that while digital retail payment platforms have replaced cash for low-ticket consumer transactions, aggregate currency demand remains elevated due to precautionary savings and store-of-value preferences in the informal economy.

\subsection{Technology Acceptance Models (UTAUT)}
The Unified Theory of Acceptance and Use of Technology (UTAUT) \cite{venkatesh2003} and its consumer extension UTAUT2 \cite{venkatesh2012} provide the principal theoretical scaffolding for technology adoption. In Indian retail payment environments, Slade et al. \cite{slade2015} and Oliveira et al. \cite{oliveira2016} demonstrate that multi-factor authentication (SMS OTPs, 4-to-6 digit UPI PINs, complex navigation) introduces significant cognitive friction for individuals with limited formal literacy. Furthermore, Talwar et al. \cite{talwar2020} and Goyal \cite{goyal2026} observe that ``transaction failure dread'' generates acute risk aversion.

\subsection{Behavioral Economics: The ``Cash-Out Loop''}
Classical adoption frameworks assume rational utility maximization based purely on speed and cost:
\begin{enumerate}
    \item \textbf{Status-Quo Bias \& Habit Persistence:} Samuelson and Zeckhauser \cite{samuelson1988} proved that decision-makers disproportionately default to established behaviors even when transition costs are objectively minimal.
    \item \textbf{Mental Accounting \& The Pain of Paying:} Prelec and Loewenstein \cite{prelec1998}, Thaler \cite{thaler1999}, and Soman \cite{soman2001} demonstrated that physical currency separation creates an immediate psychological ``pain of paying'' that serves as an automated self-regulation mechanism.
    \item \textbf{The ``Cash-Out Loop'':} In emerging markets, when informal workers or Direct Benefit Transfer (DBT) recipients receive funds into their bank accounts, their initial response is to withdraw 100\% of the cash via ATMs or Business Correspondents \cite{demirguc2022, cgap2021, sutradhar2024}. Once liquidity enters physical cash form, all downstream retail transactions inevitably occur in cash.
\end{enumerate}

\subsection{Merchant-Side Tax-Compliance Friction}
A critical institutional dynamic shaping merchant acceptance is the intersection of digital payment trails with taxation:
\begin{enumerate}
    \item \textbf{Statutory Reality vs. Rumor:} While P2P and standard bank P2M UPI transactions carry zero direct transaction tax, automated reporting to the Central Board of Direct Taxes (CBDT) under the Annual Information Statement (AIS) records all digital receipts.
    \item \textbf{GST Section 63 Threshold Enforcement:} Small retailers whose digital receipts exceed statutory GST turnover thresholds (INR~20L for services, INR~40L for goods under Section 22) face automated mismatch notices and best-judgment assessment under Section 63 for operating unregistered.
    \item \textbf{The Flight to Cash:} To suppress visible digital turnover, small informal merchants frequently disable QR displays or request cash payments. This merchant-side resistance forces even digitally capable consumers to transact in physical cash.
\end{enumerate}

\section{Data Operationalization: Global Findex}

\subsection{Data Source \& Baseline Metrics}
The empirical investigation pools authentic India microdata from two official World Bank survey waves: \textbf{Global Findex 2021} (\texttt{IND\_2021\_FINDEX\_v02\_M}, $N = 3,000$) \cite{demirguc2022} and \textbf{Global Findex 2025 / 2024 Wave} (\texttt{IND\_2024\_FINDEX\_v02\_M}, $N = 3,000$, incorporating the Digital Connectivity Tracker) \cite{worldbank2025}, yielding a combined sample of $N = 6,000$ nationally representative adults. Table \ref{tab:findex_baseline} details the verified institutional baseline metrics.

\begin{table}[htbp]
\caption{Key Empirical Baseline Data from Official Institutional Sources}
\begin{center}
\begin{tabular}{lcl}
\toprule
\textbf{Metric} & \textbf{Value} & \textbf{Primary Source} \\
\midrule
India Banknotes in Circ. (Mar 2026) & INR~41.24L Cr & RBI (2026) \cite{rbi2026} \\
Banknote Value Growth (FY25--26) & 11.9\% & RBI (2026) \cite{rbi2026} \\
Monthly UPI Volume (Oct 2026) & 25.47B & NPCI (2026) \cite{npci2026} \\
Monthly UPI Value (Oct 2026) & INR~31.53L Cr & NPCI (2026) \cite{npci2026} \\
Adult Account Ownership (2021 $\rightarrow$ 2024) & 77.5\% $\rightarrow$ 89.0\% & World Bank \cite{demirguc2022, worldbank2025} \\
Inactive Account Share (2021) & 35.0\% & World Bank \cite{demirguc2022} \\
Banked Adults Paying Only Cash (2021) & 670 Million & World Bank \cite{demirguc2022} \\
Banked Women Paying Only Cash (2021) & 350 Million & World Bank \cite{demirguc2022} \\
Digital Merchant Pay Rate (2021 $\rightarrow$ 2024) & 11.9\% $\rightarrow$ 18.3\% & Findex Microdata \\
Wage Earners Paid Cash-Only (2021 $\rightarrow$ 2024) & 39.9\% $\rightarrow$ 57.4\% & Findex Microdata \\
Stated Barrier (\texttt{fin25e4}): Habit & 77.4\% (78.6\% all) & Findex 2025 \cite{worldbank2025} \\
Stated Barrier: Merchant Cash-Only & 9.8\% (9.0\% all) & Findex 2025 \cite{worldbank2025} \\
Stated Barrier: Distrust in Digital Pay & 6.5\% (4.4\% all) & Findex 2025 \cite{worldbank2025} \\
Stated Barrier: Higher Digital Costs & 2.3\% (2.2\% all) & Findex 2025 \cite{worldbank2025} \\
\bottomrule
\end{tabular}
\label{tab:findex_baseline}
\end{center}
\end{table}

\subsection{Defining the ``Capable'' Cohort}
We isolate the paradox of non-adoption despite technical readiness by filtering:
\begin{equation}
\text{Cohort}_{\text{Capable}} = \{ i \mid \text{account}_i = 1 \land \text{mobileowner}_i = 1 \land \text{internet}_i = 1 \}
\end{equation}

Within this digitally capable cohort, the primary dependent variable captures in-store retail payment behavior:
\begin{equation}
Y_i = \begin{cases} 
1 & \text{if respondent relied exclusively on cash in store} \\ 
0 & \text{if respondent completed a digital merchant payment} 
\end{cases}
\end{equation}
where $Y_i = 1 - \text{merchantpay\_dig}_i$.

\section{Empirical Methodology}

\subsection{Survey-Weighted Logistic Regression Baseline}
As an econometric benchmark, we formulate a survey-weighted Generalized Linear Model (GLM) with a binomial distribution and logit link function:
\begin{equation}
\ln\left(\frac{P(Y_i=1)}{1 - P(Y_i=1)}\right) = \beta_0 + \sum_{j=1}^{M} \beta_j X_{ij}
\end{equation}
where normalized survey sampling weights ($w_i$) are applied to preserve representative population inferences.

\subsection{Regularized Gradient Boosted Decision Trees (XGBoost)}
The primary non-linear classifier utilizes an ensemble of $K$ additive regression trees:
\begin{equation}
\hat{y}_i = \sum_{k=1}^K f_k(\mathbf{x}_i), \quad f_k \in \mathcal{F}
\end{equation}
minimizing the regularized objective function:
\begin{equation}
\mathcal{L}(\theta) = \sum_{i=1}^n w_i \cdot \ell\left(y_i, \hat{y}_i\right) + \sum_{k=1}^K \left( \gamma T_k + \frac{1}{2}\lambda \|\mathbf{w}_k\|^2 + \alpha \|\mathbf{w}_k\|_1 \right)
\end{equation}
Class imbalance is calibrated via \texttt{scale\_pos\_weight}, and $L_1/L_2$ shrinkage penalties prevent overfitting on survey microdata \cite{chen2016}.

\subsection{TreeSHAP Game-Theoretic Interpretability}
To extract transparent, mathematically robust feature attributions, we implement TreeSHAP \cite{lundberg2020}. For a given respondent instance $x$, the Shapley value for feature $j$ is calculated as:
\begin{equation}
\phi_j(x) = \sum_{S \subseteq F \setminus \{j\}} \frac{|S|!(|F| - |S| - 1)!}{|F|!} \left[ f_x(S \cup \{j\}) - f_x(S) \right]
\end{equation}
Global barrier importance is formalized as the mean absolute Shapley value: $I_j = \frac{1}{N} \sum_{i=1}^N |\phi_j^{(i)}|$.

\section{Empirical Results \& Comparative Findings}

\subsection{Two-Wave Cohort Isolation (2021 vs. 2024)}
Across $N = 6,000$ surveyed adults, the digitally capable segment grew from 837 in 2021 (27.9\%) to 1,335 in 2024 (44.5\%), providing a pooled cohort of $N_{\text{cap}} = 2,172$ observations. Within this capable sub-population, \textbf{55.7\% remain strictly cash-reliant} (58.4\% in 2021, 54.0\% in 2024), confirming the persistence of the paradox despite massive payment infrastructure diffusion.

\subsection{Comparative Model Performance (Leak-Free 5x3 CV \& Calibrated Stack)}
To prevent simultaneous-outcome leakage, concurrent digital payment proxies (\texttt{pay\_utilities}, \texttt{fin4}, \texttt{fin14a}) were strictly excluded. Across survey-weighted 5$\times$3 Repeated Stratified K-Fold CV ($N_{\text{cap}} = 2,172$), adding domain cash-inflow intensities, debit card ownership (\texttt{fin2}), B-spline age encoding, and wave-shift interactions lifts ROC-AUC from $0.7365 \pm 0.028$ (9-var baseline; $0.7915$ without debit card) to $0.8477 \pm 0.025$ (XGBoost; PR-AUC $0.8824$, Brier $0.1571$) and $0.8470 \pm 0.022$ (4-family Stacked Ensemble + Platt Sigmoid calibration). On the held-out test set ($n = 435$), the calibrated stack achieves ROC-AUC $0.8400$ and Brier $0.1614$, while Spline-Logit with out-of-fold threshold optimization achieves PR-AUC $0.8588$ and F1 $0.7913$.

\subsection{TreeSHAP Attributions \& Econometric Odds Ratios}
TreeSHAP attributions identify upstream debit/ATM card possession (\texttt{has\_debit\_card}, $I_j = 0.9089$, OR $= 0.12, p < 0.001$), active financial engagement (\texttt{borrowed\_any} $0.3153$, \texttt{saved\_any} $0.2928$), education--income interaction ($0.3020$), and formal bank wage inflows (\texttt{wage\_bank} $0.2006$, OR $= 0.42, p < 0.001$) as primary digital catalysts. Conversely, cash wage receipt (\texttt{wage\_cash}) more than doubles retail cash-reliance odds (OR $= 2.31, p < 0.001$), female gender nearly doubles cash-reliance odds (OR $= 1.96, p = 0.010$; $I_j = 0.1572$), and the \texttt{wave\_x\_urban} interaction (OR $= 0.46, p = 0.007$) confirms significantly faster cash displacement in urban India between 2021 and 2024.

\subsection{Average Marginal Effects, Fairlie Decomposition \& 2024 Tracker Audit}
Survey-weighted Average Marginal Effects (AMEs) show urban residency flipped from insignificant in 2021 ($+4.8$ pp, $p = 0.123$) to $-6.0$ pp ($p = 0.021$) in 2024; within subgroup models, the 2024 wave reduced cash reliance by $-9.0$ pp ($p = 0.001$) in urban areas vs.\ zero significant shift in rural areas ($+3.3$ pp, $p = 0.287$). Meanwhile, bank-deposited wages (\texttt{wage\_bank}) grew from $-4.5$ pp ($p = 0.218$) in 2021 to $-19.1$ pp ($p < 10^{-9}$) in 2024, while the female penalty remained virtually constant ($+9.9$ pp in 2021; $+10.0$ pp in 2024; $+15.5$ pp for age $\ge 40$, $p = 0.0015$). Non-linear Fairlie decomposition \cite{fairlie2005} reveals that 55.25\% of the 2021--2024 weighted cash-reliance drop (60.57\% to 54.80\%) is explained by endowment expansion (debit cards, bank wages, education) and 44.75\% by behavioral shifts; for the $+14.87$ pp gender gap (66.92\% female vs.\ 52.05\% male), 50.01\% is endowment-explained and 49.99\% is unexplained household friction. Finally, auditing the 2024 Connectivity Tracker ($n = 1,335$; 5$\times$3 CV AUC $= 0.8298 \pm 0.022$) shows functional SMS literacy (\texttt{con14}) lowers cash reliance by $-8.9$ pp ($p = 0.010$; SHAP $= 0.1498$), while stated merchant refusal (\texttt{fin25e4} $= 1$, $n = 711$) has no significant demographic predictors ($p > 0.07$), indicating systemic rather than demographic-specific merchant cash insistence.

\section{Policy \& Product Interventions}
\begin{itemize}
    \item \textbf{Intercepting the Cash-Out Loop:} Policy must incentivize informal sector payroll digitization via UPI AutoPay and offer micro-cashback on essential commodities to prevent complete ATM liquidation \cite{demirguc2022}.
    \item \textbf{Harmonizing GST Cliffs with Section 44AD Safe-Harbors:} While Income Tax Section 44AD already offers an enhanced INR~3 crore limit (6\% rate vs.\ 8\% cash) when $\ge 95\%$ receipts are digital, small merchants still fear crossing the INR~40L/INR~20L GST cliff (Sections 22/63). Aligning a digital-only GST composition safe-harbor with Section 44AD removes the merchant incentive to disable QR codes.
    \item \textbf{Mitigating Scam Anxiety and Trust Deficits:} Mandating automated 60-second dispute credit reversals for failed sub-INR~500 merchant transactions \cite{talwar2020, goyal2026}.
    \item \textbf{Designing for Habit Persistence in Older Demographics:} Defaulting to zero-PIN on-device wallets (UPI Lite) and voice-activated vernacular commands (via Bhashini AI) to remove cognitive friction for individuals aged 40 and above \cite{samuelson1988}.
\end{itemize}

\section{Conclusion}
This study establishes an empirical and explainable machine learning foundation to analyze India's ``Capable-but-Cash'' paradox. By integrating macroeconomic money demand literature with micro-level Global Findex data, we demonstrate that persistent cash usage among digitally equipped adults is governed by the wage cash-out loop, demographic habit inertia, and scam anxiety, reinforced by merchant-side tax compliance frictions. Addressing India's currency demand paradox requires central banks, tax authorities, and FinTech innovators to transition from basic infrastructure expansion to targeted behavioral architectures.

\bibliographystyle{IEEEtran}
\begin{thebibliography}{00}
\bibitem{agarwal2023} A. Agarwal, ``Less - Cash Societies: A Paradoxical Roadmap Revamped through Unified Payment Systems and Electronic Currencies,'' \textit{IJSR}, vol. 12, no. 10, pp. 153--158, 2023.
\bibitem{arango2015} C. Arango, K. P. Huynh, and L. Sabetti, ``Consumer payment choice: Merchant card acceptance versus pricing incentives,'' \textit{J. Bank. Finance}, vol. 55, pp. 130--141, 2015.
\bibitem{awasthy2023} S. Awasthy, R. Misra, and S. Dhal, ``Understanding the Currency Demand Paradox in India,'' \textit{Reserve Bank of India Working Paper Series}, 2023.
\bibitem{baumol1952} W. J. Baumol, ``The Transactions Demand for Cash: An Inventory Theoretic Approach,'' \textit{Q. J. Econ.}, vol. 66, no. 4, pp. 545--556, 1952.
\bibitem{cgap2021} CGAP, \textit{Financial Inclusion and the Cash-Out Loop in Emerging Markets}, World Bank Group, 2021.
\bibitem{chen2016} T. Chen and C. Guestrin, ``XGBoost: A Scalable Tree Boosting System,'' in \textit{Proc. ACM SIGKDD}, 2016, pp. 785--794.
\bibitem{demirguc2022} A. Demirg{\"u}{\c{c}}-Kunt, L. Klapper, D. Singer, and S. Ansar, \textit{The Global Findex Database 2021}, World Bank, 2022.
\bibitem{dasilva2019} D. D'Silva et al., ``The design of digital financial infrastructure: lessons from India,'' \textit{BIS Papers}, no. 106, 2019.
\bibitem{fairlie2005} R. W. Fairlie, ``An extension of the Blinder-Oaxaca decomposition technique to logit and probit models,'' \textit{J. Econ. Soc. Meas.}, vol. 30, no. 4, pp. 305--316, 2005.
\bibitem{goodhart2020} C. A. Goodhart and J. Ashworth, ``Coronavirus and the commerce of cash: the demand for banknotes in pandemic times,'' \textit{CEPR Discussion Paper}, DP14782, 2020.
\bibitem{goyal2026} M. Goyal, ``Determinants of consumer trust and adoption barriers in UPI-based payment applications among adults in India,'' \textit{Int. J. Res. Finance Manage.}, vol. 9, no. 8, pp. 874--881, 2026.
\bibitem{kumar2025} A. Kumar and A. K. Mishra, ``Understanding India's Currency Demand Absurdity: Cash vs. Digital Payment Methods,'' \textit{IJSAT}, vol. 16, no. 1, 2025.
\bibitem{lundberg2020} S. M. Lundberg et al., ``From local explanations to global understanding with explainable AI for trees,'' \textit{Nature Mach. Intell.}, vol. 2, pp. 56--67, 2020.
\bibitem{minj2025} M. Minj, ``Shift from Cash to Cashless Economy in India: Evidence from RBI \& World Bank Data,'' \textit{JADR}, vol. 16, no. 1, 2025.
\bibitem{npci2026} NPCI, \textit{UPI Monthly Product Statistics (2024--2026)}, National Payments Corporation of India, Oct. 2026.
\bibitem{oliveira2016} T. Oliveira et al., ``Mobile payment: Understanding the determinants of customer adoption,'' \textit{Comput. Hum. Behav.}, vol. 61, pp. 404--414, 2016.
\bibitem{prelec1998} D. Prelec and G. Loewenstein, ``The Red and the Black: Mental Accounting of Savings and Debt,'' \textit{Mark. Sci.}, vol. 17, no. 1, pp. 4--28, 1998.
\bibitem{rbi2026} Reserve Bank of India, \textit{Annual Report 2025--26 \& Weekly Statistical Supplement: Currency Management}, RBI Publications, 2026.
\bibitem{samuelson1988} W. Samuelson and R. Zeckhauser, ``Status quo bias in decision making,'' \textit{J. Risk Uncertain.}, vol. 1, no. 1, pp. 7--59, 1988.
\bibitem{slade2015} E. L. Slade et al., ``Modeling consumers' adoption intentions of remote mobile payments,'' \textit{Psychol. Mark.}, vol. 32, no. 8, pp. 860--873, 2015.
\bibitem{soman2001} D. Soman, ``Effects of payment mechanism on spending behavior,'' \textit{J. Consum. Res.}, vol. 27, no. 4, pp. 460--474, 2001.
\bibitem{sutradhar2024} T. K. Sutradhar, ``Status and Determinants of Digital Transactions in India,'' \textit{Vidyasagar Univ. J. Econ.}, vol. 28, 2024.
\bibitem{talwar2020} S. Talwar, A. Dhir, and A. Khalil, ``Point of adoption and beyond: Initial trust and mobile-payment continuation intention,'' \textit{J. Retailing Consum. Serv.}, vol. 55, 2020.
\bibitem{thaler1999} R. H. Thaler, ``Mental accounting matters,'' \textit{J. Behav. Decis. Making}, vol. 12, no. 3, pp. 183--206, 1999.
\bibitem{tobin1956} J. Tobin, ``The Interest-Elasticity of Transactions Demand For Cash,'' \textit{Rev. Econ. Stat.}, vol. 38, no. 3, pp. 241--247, 1956.
\bibitem{venkatesh2003} V. Venkatesh et al., ``User acceptance of information technology: Toward a unified view,'' \textit{MIS Q.}, vol. 27, no. 3, pp. 425--478, 2003.
\bibitem{venkatesh2012} V. Venkatesh, J. Y. Thong, and X. Xu, ``Consumer acceptance and use of information technology: Extending UTAUT,'' \textit{MIS Q.}, vol. 36, no. 1, pp. 157--178, 2012.
\bibitem{worldbank2023} World Bank, \textit{G20 Policy Recommendations for Advancing Financial Inclusion through Digital Public Infrastructure}, World Bank Group, 2023.
\bibitem{worldbank2025} World Bank, \textit{The Global Findex Database 2025: Connectivity and Financial Inclusion in the Digital Economy}, World Bank Group, 2025.
\end{thebibliography}

\end{document}
